\section{Tubular coordinates around the toroidal electron}
\label{app:tubular-electron}

The qualitative Williamson--van der Mark picture --- a localized
photon-like internal mode circulating at (essentially) the speed of
light along a double-loop path confined to a torus --- can be made
more explicit by introducing tubular coordinates around the soliton's
centerline. This is the continuum analogue of the ``curved tunnel''
or ``tape'' picture discussed in Section~\ref{conceptual-model}, and
will be useful for analysing how geometric nonlinearity couples the
internal Compton mode to lateral deformations of the brane. It is
important to emphasize that the ``tube'' or ``tunnel'' terminology
refers to a \emph{tubular neighbourhood} (coordinate patch) on the
brane---the region where the amplitude $|\xi|$ is appreciable. It is
not a separate material structure or an additional object with walls;
rather, it is simply a convenient way to parametrize the portion of
the single brane where the electron's internal mode is localized.

Let $\mathbf{C}(z)\in\mathbb R^3$ denote the spatial centerline of the
toroidal electron in the brane's three lateral directions
$X^1,X^2,X^3$, parametrized by arclength $z\in[0,L)$ in the electron
rest frame, so that $|\partial_z\mathbf C(z)|=1$ and $\mathbf C(z+L)
=\mathbf C(z)$. Along this closed curve we introduce the Frenet--Serret
frame $(\mathbf t(z),\mathbf n(z),\mathbf b(z))$ with curvature
$\kappa(z)$ and torsion $\tau(z)$,
%
\begin{equation}
  \mathbf t'=\kappa\,\mathbf n,\qquad
  \mathbf n'=-\kappa\,\mathbf t+\tau\,\mathbf b,\qquad
  \mathbf b'=-\tau\,\mathbf n,
  \label{eq:frenet-serret}
\end{equation}
%
where $'=\partial_z$. A small tubular neighbourhood of the centerline
is then parametrized by coordinates $(z,x,y)$ via
%
\begin{equation}
  \mathbf r(z,x,y)
  \;=\;
  \mathbf C(z) + x\,\mathbf n(z) + y\,\mathbf b(z),
  \label{eq:tubular-embedding}
\end{equation}
%
with $(x,y)$ spanning the cross-section of the ``tunnel'' around the
electron. The microscopic embedding field is still
$\mathbf X=(X^1,X^2,X^3,X^4)$, but in these adapted coordinates the
three lateral components $(X^1,X^2,X^3)$ are replaced by
$\mathbf r(z,x,y)$, while the amplitude direction $X^4$ remains a
separate normal coordinate. We denote the microscopic amplitude mode in
this tubular neighbourhood by
%
\begin{equation}
  \xi(t,z,x,y) \;:=\; X^4\bigl(t,\mathbf r(z,x,y)\bigr),
  \label{eq:xi-from-X4}
\end{equation}
%
so that $(t,z,x,y)$ serve as local coordinates on the world-tube of
the electron---that is, the spacetime region of the brane where the
amplitude $|\xi|$ is appreciable---with the double-loop structure
encoded in the internal mode pattern of $\xi$ across the cross-section
$(x,y)$.

The induced metric $g_{ab}$ on the brane in the coordinates
$(z,x,y)$ is obtained by pulling back the flat ambient metric
in the lateral directions:
%
\begin{equation}
  g_{ab}(z,x,y)
  \;=\;
  \partial_a\mathbf r\cdot\partial_b\mathbf r,
  \qquad a,b\in\{z,x,y\}.
  \label{eq:tubular-metric-def}
\end{equation}
%
Using the Frenet--Serret relations~\eqref{eq:frenet-serret} one finds
to first order in the transverse distance $\rho=\sqrt{x^2+y^2}$ from
the centerline,
%
\begin{equation}
  g_{ab}
  \;\approx\;
  \begin{pmatrix}
    1 - 2\kappa x & -\,\tau y & \;\;\tau x \\[0.4ex]
    -\,\tau y & 1 & 0 \\[0.4ex]
    \;\;\tau x & 0 & 1
  \end{pmatrix},
  \qquad
  g^{ab}
  \;\approx\;
  \begin{pmatrix}
    1 + 2\kappa x & \;\;\tau y & -\,\tau x \\[0.4ex]
    \;\;\tau y & 1 & 0 \\[0.4ex]
    -\,\tau x & 0 & 1
  \end{pmatrix},
  \label{eq:tubular-metric}
\end{equation}
%
where $g^{ab}$ is the inverse metric, and we have neglected terms of
order $\kappa^2\rho^2$ and $\tau^2\rho^2$. In these adapted
coordinates the toroidal soliton appears, locally, as a straight tunnel
aligned with the $z$-axis; the global curvature and double-loop
structure are encoded entirely in the $z$-dependence of the frame
$(\mathbf n,\mathbf b)$ and the metric coefficients $g_{ab}(z,x,y)$.
The electromagnetic energy circulating along the double loop is then
represented by a structured internal mode $\xi(t,z,x,y)$ localized
within this world-tube (the brane region where $|\xi|$ is large),
confined to a finite range of $(x,y)$ around the centerline.

\paragraph{Effective amplitude dynamics in tubular coordinates.}

Restricting attention to the dominant microscopic amplitude mode
$\xi(t,z,x,y)$ inside the electron's world-tube, we can write a
minimal effective Lagrangian density in the adapted coordinates
$(t,z,x,y)$ of the form
%
\begin{equation}
  \mathcal{L}_\xi
  \;=\;
  \frac{\rho_m}{2}\,(\partial_t \xi)^2
  \;-\;
  \frac{T_\mathrm{eff}}{2}\,g^{ab}\,\partial_a\xi\,\partial_b\xi
  \;-\;
  V_\mathrm{nl}(\xi,\partial\xi,\dots),
  \label{eq:xi-tubular-lagrangian}
\end{equation}
%
where $\rho_m$ is the effective mass density introduced in
Section~\ref{sec:emergent-relativity}, $T_\mathrm{eff}$ is an effective
tension along the tube that depends on the microscopic spring constant
and rest length, and $V_\mathrm{nl}$ collects residual nonlinear
corrections beyond the geometric contributions contained in the metric
$g^{ab}$. Inserting the approximate inverse metric
\eqref{eq:tubular-metric} yields, to first order in curvature and
torsion,
%
\begin{align}
  g^{ab}\,\partial_a\xi\,\partial_b\xi
  \;\approx\;
  &\bigl(1 + 2\kappa x\bigr)\,(\partial_z\xi)^2
  + (\partial_x\xi)^2
  + (\partial_y\xi)^2
  \nonumber\\[0.5ex]
  &\quad
  + 2\tau y\,(\partial_z\xi)(\partial_x\xi)
  - 2\tau x\,(\partial_z\xi)(\partial_y\xi),
  \label{eq:xi-gradient-tubular}
\end{align}
%
so that the gradient energy density of the amplitude mode takes the
schematic form
%
\begin{equation}
  \mathcal{E}_{\mathrm{grad}}
  \;\approx\;
  \frac{T_\mathrm{eff}}{2}\Bigl[
    \bigl(1+2\kappa x\bigr)(\partial_z\xi)^2
    + (\partial_x\xi)^2
    + (\partial_y\xi)^2
    + 2\tau y\,(\partial_z\xi)(\partial_x\xi)
    - 2\tau x\,(\partial_z\xi)(\partial_y\xi)
  \Bigr].
  \label{eq:xi-gradient-energy}
\end{equation}
%
Two features are worth emphasizing. First, the curvature term
$\propto\kappa x(\partial_z\xi)^2$ implies that for sufficiently
large longitudinal gradients of $\xi$ it becomes energetically
favourable for the mode to shift its support away from the
centerline in the transverse directions $(x,y)$, i.e.\ to develop a
lateral ``bulge''. Second, the torsion terms
$\propto\tau\,(\partial_z\xi)(\partial_x\xi)$ and
$\propto\tau\,(\partial_z\xi)(\partial_y\xi)$ show that a purely
longitudinal oscillation is not an eigenmode once the tube is twisted:
steep variations along $z$ inevitably source transverse gradients
$\partial_x\xi$ and $\partial_y\xi$. Together, these geometric
couplings provide a concrete mechanism by which the strongly excited,
Compton-frequency internal mode of the toroidal electron drives lateral
deformations of the brane that confine the energy to a finite tubular
region, while the microscopic spring rest length controls the relative
cost of longitudinal stretching versus transverse bulging through the
effective tension $T_\mathrm{eff}$ and thus sets the critical amplitude
threshold for self-confinement.

In the brane picture, this localized, circulating electromagnetic energy
density produces a quasi-static amplitude deformation $\bar{X}^4$ in
its vicinity. At the coarse-grained level we can thus treat the
effective charge density $\rho(\mathbf{x})$ as a functional of the
time-averaged internal energy distribution,
\begin{equation}
  \rho(\mathbf{x})
  \;\propto\;
  \chi\,\big\langle u_{\text{EM}}(\mathbf{x},t)\big\rangle_{\text{fast}},
  \label{eq:rho-from-u}
\end{equation}
where $u_{\text{EM}}$ denotes the electromagnetic energy density of the
confined toroidal mode and $\chi=\pm 1$ encodes the orientation
(chirality) of the internal field configuration. The sign of $\chi$
distinguishes electron and positron: reversing the direction of the
internal circulation flips the sign of the induced amplitude deformation
and hence of $\Phi(\mathbf{x})$, while leaving the absolute charge
magnitude and magnetic dipole moment unchanged, in agreement with
experiment.

Far outside the Compton-scale core, the rapid internal precession of
the toroidal configuration renders the time-averaged deformation
$\bar{X}^4(\mathbf{x})$ effectively spherically symmetric. In this
regime the scalar potential defined in Eq.~\eqref{eq:phi-from-xi4}
obeys a Poisson-type equation on the brane,
\begin{equation}
  \nabla^2 \Phi(\mathbf{x})
  \;=\;
  -\,\frac{\rho_{\text{free}}(\mathbf{x})}{\varepsilon_{\mathrm{eff}}},
  \label{eq:poisson-from-brane}
\end{equation}
with an effective permittivity $\varepsilon_{\mathrm{eff}}$ determined
by the elastic response of the medium. Here $\rho_{\text{free}}$ denotes
the net (monopole) free charge density: bound, polarization-like
patterns of $X^4$ that cancel on scales larger than the Compton
wavelength do not contribute to this source term.

\paragraph{Effective Poisson closure.}
Equation~\eqref{eq:poisson-from-brane} is, at this stage, an \emph{effective closure}: we impose a phenomenological constitutive law that links the coarse-grained amplitude $\bar X^4$ to an electrostatic potential $\Phi_{\text{EM}}$ satisfying a Poisson equation. This closure is chosen such that localized brane deformations give rise to Coulomb-like fields at large distances, in agreement with observed electrostatic behavior. A microscopic derivation of \eqref{eq:poisson-from-brane} directly from the brane Lagrangian $L[\mathbf X]$---for instance via a systematic multipole analysis of static solutions---is left as future work.

Combining Eqs.~\eqref{eq:phi-from-xi4} and~\eqref{eq:poisson-from-brane}
yields the local geometric relation
\begin{equation}
  \rho_{\text{free}}(\mathbf{x})
  = -\,\varepsilon_{\mathrm{eff}}\kappa\,\nabla^2 \bar{X}^4(\mathbf{x}),
\end{equation}
so that net charge is associated not with the absolute value of the
amplitude displacement, but with its curvature pattern on the brane.

The resulting field
\[
  \mathbf{E}(\mathbf{x}) = -\nabla\Phi(\mathbf{x})
  = -\,\kappa\,\nabla \bar{X}^4(\mathbf{x})
\]
reproduces the familiar Coulomb behavior of a point charge at large
distances, even though microscopically the source is an extended,
toroidal soliton of finite size. In the static rest frame of the
soliton, and for this effective free charge component, the electric
field is literally proportional to the gradient of the brane surface in
the fourth (amplitude) dimension. From this perspective, regions of
positive and negative free charge correspond to peaks and troughs of the
brane surface in the amplitude direction, with the sign of the gradient
fixing the sign of the observed electric field.

It is important to emphasize that the same amplitude deformation that
encodes electric charge also participates in the emergent gravitational
response of the brane. However, these are distinct effects arising from
different aspects of the deformation field $X^4$:

\begin{itemize}
\item \textbf{Gravitational sector:} responds to the induced metric
  $g_{ij}[\mathbf{X}]$, which is roughly quadratic in gradients of $X^4$.
  The isotropic component of the induced strain pattern---set by the total
  internal energy of the toroidal mode---contributes to curvature and
  inertial mass.

\item \textbf{Electromagnetic sector:} couples directly to the time-averaged
  amplitude $\bar{X}^4$ and its Laplacian via
  Eqs.~\eqref{eq:phi-from-xi4} and~\eqref{eq:poisson-from-brane}. The signed
  component associated with the orientation $\chi$ encodes the electric
  interaction in the rest frame.
\end{itemize}

Thus, the same underlying amplitude deformation yields two distinct
effective fields. In this way, a single localized brane deformation
unifies mass, charge, and the external electromagnetic potential as
different macroscopic manifestations of the same microscopic structure,
without double-counting: gravity depends on metric curvature (quadratic
in $\nabla X^4$), while electrostatics depends linearly on the Laplacian
of $\bar{X}^4$.